
          %%%%% ~~~~~~~~~~~~~~~~~~~~ %%%%%

\section{Target period, land cover typology and target variables}

          %%%%% ~~~~~~~~~~~~~~~~~~~~ %%%%%

\subsection{Target period}
The target period of this release of land cover area estimates is circa 2020. The slight 
temporal imprecision of the target period is due to the limited availability of high-resolution 
satellite imagery data (needed to generate reference data as input to the estimation 
procedure). 
 
\subsection{Target population and sampled population}
The target population was the collection of 100 metre $\times$ 100 metre tiles (grid cells) 
overlapping the land mass of Canada from the CoE reference geospatial grid of spatial 
resolution of 100 metres $\times$ 100 metres. The sampled population was identical to the target 
population; in other words, there were no coverage errors. As per standard nomenclature 
from finite-population sampling, we will also refer to the grid cells as sampling units in what 
follows. 
 
\subsection{Simplified Land Cover Register (LCR) map typology}
The land cover typology used for this data product is a slightly simplified version of the LCR 
map typology. This simplified typology is obtained from the LCR map typology by collapsing 
the classes treed (non-wetland) and treed wetland. This simplified typology was used due to 
the infeasibility of distinguishing treed uplands from treed wetlands based on satellite 
imagery. In what follows, we will refer to this typology as the simplified LCR map typology. 
 
\subsection{Reference land cover type versus stratification land cover type}
While we used only one land cover typology (i.e., the simplified RLC map typology), we used 
this typology for two distinct purposes: (i) defining the target variables (hence also 
generating reference data), and (ii) stratification of the target population in order to create 
the sampling design (see below). Both usages appear in the estimation procedure, and they 
must be kept conceptually distinct to understand and execute the estimation procedure. In 
what follows, we will use the term
\textit{reference land cover type} (or, more briefly, \textit{reference class})
in relation to purpose (i) and
\textit{stratification land cover type} (or, more briefly, \textit{map class})
in relation to purpose (ii). 
 
\subsection{Target variables}
The target variables are the areas (in square kilometres) around the target period (calendar 
year 2020) of the different reference land cover types within each target geography. The 
target geographies are ecoprovinces, ecozones, sub-drainage areas, major drainage areas, 
drainage regions, and all of Canada. Note that the ecoprovinces are nested within ecozones, 
while the sub-drainage areas are nested within the major drainage areas.

          %%%%% ~~~~~~~~~~~~~~~~~~~~ %%%%%
